\section{Physics Geometry: Coordinate and Vector Conventions}

\subsection{Project-Wide Rule}

No array or column may be named only \texttt{x}, \texttt{y}, \texttt{vx}, or
\texttt{vy} at a boundary between physics modules unless its coordinate frame
and units are defined by the enclosing typed object. Points and vectors must be
transformed separately.

\subsection{Coordinate Frames}

\paragraph{Image frame.}
The image frame uses \texttt{x\_px} and \texttt{y\_px}. The origin is the
top-left video pixel, \(+x\) points right, \(+y\) points downward, and units are
pixels.

\paragraph{Device Cartesian frame.}
The device frame uses \texttt{x\_device\_um} and \texttt{y\_device\_um}. The
origin is fixed by the lower-left image reference line, \(+x\) points right,
\(+y\) points upward, and units are micrometers. For \texttt{video\_2}, the
fixed reference is taken from the device region metadata image height:
\[
  y_{\mathrm{reference,px}} = 596.
\]

\paragraph{Frozen CFD native frame.}
The frozen CFD Version 1 mesh uses \texttt{x\_cfd\_um} and
\texttt{y\_cfd\_um}. This frame inherits the calibrated full-image centerline:
\(+x\) points right, \(+y\) points downward, and units are micrometers. This is
the native frame of the stored CFD coefficients; it is not the preferred
cross-module Cartesian vector frame.

\subsection{Point Transforms}

For pixel scale \(s\) in micrometers per pixel,
\[
  x_{\mathrm{device}} = s x_{\mathrm{px}},
  \qquad
  y_{\mathrm{device}} = s(y_{\mathrm{reference,px}} - y_{\mathrm{px}}).
\]

The frozen CFD mesh uses the calibrated image-oriented coordinate system, so
\[
  x_{\mathrm{cfd}} = x_{\mathrm{device}},
  \qquad
  y_{\mathrm{cfd}} =
  y_{\mathrm{reference,um}} - y_{\mathrm{device}}.
\]

\subsection{Vector Transforms}

Vectors use only the linear part of the affine transform. Translation is never
applied to a vector:
\[
  v_{x,\mathrm{device}} = s v_{x,\mathrm{image}},
  \qquad
  v_{y,\mathrm{device}} = -s v_{y,\mathrm{image}}.
\]

Between device Cartesian and CFD native frames:
\[
  u_{x,\mathrm{cfd}} = u_{x,\mathrm{device}},
  \qquad
  u_{y,\mathrm{cfd}} = -u_{y,\mathrm{device}}.
\]

The inverse vector transform has the same y reflection:
\[
  u_{x,\mathrm{device}} = u_{x,\mathrm{cfd}},
  \qquad
  u_{y,\mathrm{device}} = -u_{y,\mathrm{cfd}}.
\]

\subsection{Observed Droplet Velocities}

The current tracked-feature file contains framewise centroid positions and no
canonical physical velocity columns. Centered finite differences of
\texttt{centroid\_x} and \texttt{centroid\_y} produce image-frame velocities in
pixels per frame. The device-frame conversion is:
\[
  v_{x,\mathrm{device}} = s v_{x,\mathrm{image}},
  \qquad
  v_{y,\mathrm{device}} = -s v_{y,\mathrm{image}}.
\]

No frame-rate conversion is currently applied in the enrichment direction audit,
so observed droplet speeds remain in micrometers per frame. Direction
comparisons are still valid because both observed and CFD vectors are converted
to the same device Cartesian frame before computing angles.

\subsection{CFD Direction Audit}

The coordinate audit evaluates the frozen CFD Version 1 field directly at
representative inlet and outlet points. For a point \(\mathbf{x}\), junction
point \(\mathbf{x}_j\), and velocity \(\mathbf{u}\), the geometric tests are:
\[
  \mathbf{u}\cdot(\mathbf{x}_j-\mathbf{x}) > 0
  \quad\mathrm{in~the~inlet},
\]
\[
  \mathbf{u}\cdot(\mathbf{x}-\mathbf{x}_j) > 0
  \quad\mathrm{in~each~outlet}.
\]

\begin{center}
\begin{tabular}{lrrrrr}
\toprule
Region & \(x_{\mathrm{cfd}}\) & \(y_{\mathrm{cfd}}\) & \(u_x\) & \(u_y\) & Pass \\
\midrule
inlet & 1320.000 & 590.000 & \(-3.115{\times}10^{-9}\) & \(8.167{\times}10^{-2}\) & true \\
left outlet & 1046.000 & 860.000 & \(-4.021{\times}10^{-2}\) & \(-8.003{\times}10^{-5}\) & true \\
right outlet & 1593.934 & 861.612 & \(4.050{\times}10^{-2}\) & \(1.281{\times}10^{-3}\) & true \\
\bottomrule
\end{tabular}
\end{center}

The signed boundary fluxes are:
\[
  Q_{\mathrm{inlet}}=-5.444444{\times}10^{-6}~\mathrm{m^2/s},
  \qquad
  Q_{\mathrm{left}}=Q_{\mathrm{right}}=2.722222{\times}10^{-6}~\mathrm{m^2/s}.
\]

The negative inlet sign follows the outward-normal convention: flow enters the
domain at the inlet and leaves through the outlets. The frozen CFD field is
therefore physically oriented correctly. The misleading inlet arrow was a
downstream convention problem: diagnostics and enrichment columns previously
used ambiguous y-down quantities without naming the frame explicitly.

\begin{figure}[H]
  \centering
  \includegraphics[width=0.90\linewidth,height=0.65\textheight,keepaspectratio]{../outputs/physics/coordinate_audit/figures/cfd_native_direction_check.png}
  \caption{Direction check in the frozen CFD native frame. This frame has \(+y\) downward.}
  \label{fig:coordinate-audit-cfd}
\end{figure}

\begin{figure}[H]
  \centering
  \includegraphics[width=0.90\linewidth,height=0.65\textheight,keepaspectratio]{../outputs/physics/coordinate_audit/figures/device_cartesian_direction_check.png}
  \caption{Direction check in the device Cartesian frame. This frame has \(+y\) upward and uses no axis inversion.}
  \label{fig:coordinate-audit-device}
\end{figure}

\begin{figure}[H]
  \centering
  \includegraphics[width=0.90\linewidth,height=0.65\textheight,keepaspectratio]{../outputs/physics/coordinate_audit/figures/image_coordinate_direction_check.png}
  \caption{Direction check in image coordinates. The inlet arrow points visually from the upper stem toward the junction.}
  \label{fig:coordinate-audit-image}
\end{figure}
